%  MLE-AI-201 Machine Learning -- Final Project report
%  Self-contained IEEE-style two-column layout (standard packages only).
%
%  Build (from report/):   pdflatex report.tex ; pdflatex report.tex
%  Every table in tables/ and figure in figures/ is generated by
%      python -m src.run_all
\documentclass[10pt,twocolumn,a4paper]{article}

\usepackage[a4paper,top=0.75in,bottom=1.0in,left=0.68in,right=0.68in,columnsep=0.25in]{geometry}
\usepackage{mathptmx}
\usepackage[T1]{fontenc}
\usepackage[utf8]{inputenc}
\usepackage{graphicx}
\usepackage{amsmath,amssymb}
\usepackage{booktabs}
\usepackage{array}
\usepackage{titlesec}
\usepackage{caption}
\usepackage{cite}
\usepackage[hidelinks]{hyperref}
\usepackage{url}

\renewcommand{\thesection}{\Roman{section}}
\renewcommand{\thesubsection}{\Alph{subsection}}
\titleformat{\section}[block]{\centering\normalsize\scshape}{\thesection.}{0.5em}{}
\titleformat{\subsection}[block]{\normalsize\itshape}{\thesubsection.}{0.5em}{}
\titlespacing*{\section}{0pt}{1.4ex plus 1ex minus .2ex}{1ex}
\titlespacing*{\subsection}{0pt}{1.0ex plus 1ex minus .2ex}{0.6ex}
\newcommand{\abstractkw}[1]{\noindent\textbf{\textit{#1}}}
\captionsetup{font=small,labelsep=period}
\setlength{\parindent}{1em}
\sloppy
\graphicspath{{figures/}}

\newcommand{\repourl}{github.com/pnah6699-maker/bina-az}

\begin{document}

\twocolumn[
  \begin{@twocolumnfalse}
  \begin{center}
    {\LARGE\bfseries Pricing Baku Homes from Scratch:\\
     A NumPy Decision Tree and Pegasos SVM on bina.az Listings\par}
    \vspace{1.2em}
    {\large
      Farhad Budagov\quad Murad Atlukhanov\quad Panah Gulueyev\quad Mikayil Hasanov\par}
    \vspace{0.5em}
    {\itshape MLE-AI-201 Machine Learning, AI Academy --- Cohort II\par}
    \vspace{1.4em}
  \end{center}
  \begin{quote}
    \small
    \abstractkw{Abstract---}%
    We predict residential sale prices in Azerbaijan from a raw scrape of
    100{,}775 bina.az listings. After de-duplicating repeat scrapes,
    removing the target-leaking columns \texttt{unit\_price} and
    \texttt{total\_price}, and applying fixed domain outlier rules,
    57{,}473 listings with 92 features remain. We frame two tasks: (A)
    regression of $\log(\text{price})$ and (B) a premium/standard tier
    split at the training-set median of 209{,}000~AZN. We implement a CART
    decision tree (Gini, entropy and MSE criteria with four stopping rules)
    and a soft-margin SVM trained with the Pegasos sub-gradient method,
    extended with random Fourier features, using NumPy only. On a held-out
    test set our tree reaches RMSE $0.232$ on $\log(\text{price})$
    ($R^2=0.864$, MAPE 15.4\%) and F1 $0.921$ on the tier task. The linear
    Pegasos SVM reaches F1 $0.900$ with ROC-AUC $0.965$. Each result is
    within $0.001$ of the matched scikit-learn baseline and within one
    cross-validation standard deviation. A random forest built from our own
    trees lowers the regression error to $0.185$. The errors concentrate on
    the cheapest and most expensive listings, houses and villas, and
    sparsely listed micro-locations.
    Code: \href{https://\repourl}{\texttt{\repourl}} (tag \texttt{v1.0-final}).

    \vspace{0.6em}
    \abstractkw{Index Terms---}decision tree, support-vector machine,
    Pegasos, random Fourier features, regression, classification,
    real-estate price prediction.
  \end{quote}
  \vspace{1.0em}
  \end{@twocolumnfalse}
]

\section{Introduction}
A listing price on bina.az, Azerbaijan's largest real-estate portal, is the
seller's opening position in a negotiation. An estimate built from a home's
size, location and condition helps buyers spot over-priced listings and
helps sellers price realistically. The data is a genuine scrape, so most of
the work is turning it into an honest modelling dataset. Our contributions:
\begin{itemize}\setlength\itemsep{0pt}
  \item a documented cleaning pipeline that removes repeat scrapes,
    target leakage and data-entry errors, with every row count logged
    (Table~\ref{tab:clean});
  \item a from-scratch CART tree whose splits provably match scikit-learn
    on a fully determined toy example, sped up by pre-sorting each feature
    once;
  \item a from-scratch Pegasos SVM with margin and support-vector
    reporting, plus a random-Fourier-feature (RFF) non-linear version;
  \item a leakage-free protocol (validation-only tuning, re-fitting all
    preprocessing inside every CV fold) and a matched benchmark against
    scikit-learn, followed by an error analysis;
  \item all three bonuses: Ridge via the normal equations, a random forest
    of our own trees, and k-means/PCA geospatial analysis.
\end{itemize}

\section{Background}
\textit{Decision trees} split a node with $N$ samples on the
feature/threshold pair that maximises the information gain
$\mathrm{IG}=I(\text{parent})-\sum_{c}\tfrac{N_c}{N}I(c)$. Here $I$ is the
Gini impurity $1-\sum_k p_k^2$, the entropy $-\sum_k p_k\log_2 p_k$, or,
for regression, the node variance
$\tfrac1N\sum_i(y_i-\bar y)^2$~\cite{cart}. \textit{Soft-margin SVMs}
minimise the regularised hinge loss
\begin{equation}
  J(\mathbf{w},b)=\tfrac{\lambda}{2}\lVert\mathbf{w}\rVert^{2}
  +\tfrac{1}{n}\sum_{i=1}^{n}\max\!\big(0,\,1-y_i(\mathbf{w}^{\top}\mathbf{x}_i+b)\big),
  \label{eq:svm}
\end{equation}
which is equivalent to the usual $C$-formulation with $C=1/(\lambda n)$.
\textit{Pegasos}~\cite{pegasos} minimises (\ref{eq:svm}) by stochastic
sub-gradient steps with $\eta_t=1/(\lambda t)$ and needs no QP or SMO
solver. \textit{Random Fourier features}~\cite{rff} map $\mathbf{x}$ to
$\mathbf{z}(\mathbf{x})=\sqrt{2/D}\cos(W\mathbf{x}+\mathbf{c})$ with
$W_{ij}\sim\mathcal N(0,2\gamma)$, so that
$\mathbf z(\mathbf x)^{\top}\mathbf z(\mathbf x')\approx
e^{-\gamma\lVert\mathbf x-\mathbf x'\rVert^2}$. A linear SVM on
$\mathbf z$ then approximates an RBF-kernel SVM.

\section{Data and Preprocessing}
\subsection{Dataset and cleaning}
We use the Kaggle dataset\\ \url{sehriyarmemmedli/binaaz-sale-project}
(version~1, file \texttt{house\_sale.csv}): 100{,}775 rows and 51
columns, scraped 5--16 October 2024. Every price is quoted in AZN; the code
still maps currencies explicitly. Table~\ref{tab:clean} lists each
cleaning step.
(i)~\emph{Duplicates:} 36{,}321 rows repeat an item URL scraped on
another day, so we keep the latest scrape per URL. A further 753 re-posts
share price, area, rooms, floor, coordinates and category.
(ii)~\emph{Scope:} we keep the three residential categories: new
builds, old stock and houses/villas. Land is priced per \textit{sot}
(100\,m$^2$), and commercial, office and garage listings form other
markets.
(iii)~\emph{Type coercion:} strings such as \texttt{"145 m²"},
\texttt{"4.5 sot"} and \texttt{"7 / 9"} become area in m$^2$, land area in m$^2$,
floor and total floors. The district (\textit{rayon}) and a metro tag are
parsed from \texttt{extra\_info}.
(iv)~\emph{Outliers}, with fixed domain rules rather than data-driven
percentiles, so that no statistic is computed on test rows: price in
[10\,k, 5\,M]~AZN, area in [15, 1000]\,m$^2$, a price per m$^2$ in
[300, 15\,000] used only to drop entry errors and never as a feature, and
coordinates inside Azerbaijan. Rooms above 10 and inconsistent floors
(floor $>$ total, or more than 40 floors) are set to missing rather than
dropped.
(v)~\emph{Missing values} (rooms 9\%, floor for apartments, repair
status 6\%) are imputed with the training median, and a missing-indicator
column is added for each. Houses have no floor, so their indicator is
always 1.

\begin{table}[t]
  \centering\small
  \caption{Cleaning log (rows left after each step).}
  \label{tab:clean}
  \begin{tabular}{@{}p{0.70\linewidth}r@{}}
\toprule
Step & Rows left \\
\midrule
    raw rows & 100,775 \\
    known currency, numeric price & 100,775 \\
    unique listing URL (latest scrape) & 64,454 \\
    residential categories & 58,452 \\
    no exact re-posts (same price/area/rooms/floor/coords) & 57,699 \\
    price in [10,000, 5,000,000] AZN & 57,682 \\
    area in [15, 1000] m$^2$ & 57,594 \\
    price/m$^2$ in [300, 15,000] (entry errors) & 57,474 \\
    coordinates inside Azerbaijan & 57,473 \\
\bottomrule
\end{tabular}

\end{table}

\subsection{Leakage}
\texttt{unit\_price} (``3\,450 AZN/m²'') equals price divided by area, and
\texttt{total\_price} equals \texttt{price} in all but 13 rows. Both are
arithmetic functions of the target and are dropped before any other step
(\texttt{LEAKAGE\_COLUMNS} in \texttt{src/data\_prep.py}). We also drop
identifiers and personal data (\texttt{owner\_name}, \texttt{shop\_name},
raw \texttt{address}, URLs and image links), scrape metadata, free text,
and \texttt{views}. View counts accumulate \emph{after} a listing is posted
and partly react to its price, so they would not be available at
prediction time. A unit test (\texttt{test\_leakage\_columns\_are\_dropped})
checks that no feature name contains ``price''.

\subsection{Features, tasks and split}
The 92 features are: 13 numeric (area, $\log$ area, rooms, area per room,
floor, total floors, floor ratio, first-floor and top-floor flags, land
area, lat, lng, and haversine distance to the city centre); 9 binary
(repaired, bill of sale, mortgage, agent, residential complex, VIP,
featured, Baku, metro tag); 4 missing-value flags; and one-hot columns for
category, 11 districts and 46 micro-locations. Districts and locations
with fewer than 150 training rows are pooled into ``other''. We use one-hot
encoding because it imposes no false order on districts and suits both
trees and the SVM. Every learned quantity in the \texttt{Featurizer}
(medians and category levels) is fit on the training split only.

\emph{Task A} regresses $\log(\text{price})$, since the raw price has
skewness 5.5 against 0.29 after the log (Fig.~\ref{fig:eda}a).
\emph{Task B} labels a listing \emph{premium} if its price exceeds the
\emph{training} median of 209{,}000~AZN. The validation and test labels
reuse that threshold. The 70/15/15 split is stratified on price deciles,
so all three splits have the same price distribution and therefore the
same tier balance (Table~\ref{tab:data}).

\begin{table}[t]
  \centering\small
  \caption{Dataset summary after cleaning.}
  \label{tab:data}
  \begin{tabular}{@{}lr@{}}
\toprule
Property & Value \\
\midrule
Rows (after cleaning) & 57,473 \\
Features used & 92 \\
Train / Val / Test & 40,233 / 8,620 / 8,620 \\
Tier threshold (train median) & 209,000 AZN \\
Premium share train/val/test (\%) & 49.7 / 49.8 / 49.6 \\
\bottomrule
\end{tabular}

\end{table}

\begin{figure*}[t]
  \centering
  \includegraphics[width=0.56\linewidth]{eda_price_distribution.pdf}\hfill
  \includegraphics[width=0.40\linewidth]{eda_location.pdf}\\[2pt]
  \includegraphics[width=0.31\linewidth]{eda_area_vs_price.pdf}\hfill
  \includegraphics[width=0.31\linewidth]{eda_rooms.pdf}\hfill
  \includegraphics[width=0.31\linewidth]{eda_district.pdf}
  \caption{EDA on the 57{,}473 cleaned listings. (a)~Price is heavily
  right-skewed and $\log(\text{price})$ is nearly symmetric; the dashed
  line is the tier threshold. (b)~Price per m$^2$ peaks in central Baku
  (Sabail, Nasimi, Yasamal). (c)~Price scales roughly as a power of
  area, with houses priced lower per m$^2$. (d)~Price grows
  monotonically with rooms. (e)~The median price differs about 4$\times$
  between Sabail and the outer districts.}
  \label{fig:eda}
\end{figure*}

\section{Method}
\subsection{Decision tree (from scratch)}
\texttt{DecisionTree(task, criterion, max\_depth, min\_samples\_split,
min\_samples\_leaf, min\_impurity\_decrease)} mirrors scikit-learn's API.
\emph{Split search:} each feature is sorted once at the root. A child
inherits its rows in sorted order by stably filtering its parent's order,
which costs $O(Nd)$ per level instead of $O(Nd\log N)$. For a node, all
$N-1$ cut positions of all features are scored in one vectorised pass.
Classification uses cumulative class counts, giving $p_k$ on each side of
every cut. Regression uses cumulative $\sum y$ and $\sum y^2$ on centred
targets, giving the SSE on each side. A cut is valid only between two
distinct values and only if both children keep at least
\texttt{min\_samples\_leaf} rows. The threshold is the midpoint of the two
neighbouring values, and a row goes left iff $x_j\le t$.
\emph{Stopping:} a node becomes a leaf at \texttt{max\_depth}, when it has
fewer than \texttt{min\_samples\_split} rows, when it is pure, or when the
weighted decrease $\frac{N_t}{N}\,\mathrm{IG}$ falls below
\texttt{min\_impurity\_decrease} (scikit-learn's definition). A unit test
shows that each of the four rules changes the number of leaves.
\emph{Leaves} predict the mean (regression) or the class frequencies,
which give \texttt{predict\_proba}, with the majority class as the label.
\emph{Importances} sum each feature's weighted impurity decrease over all
splits; they match scikit-learn to $10^{-10}$ in a test. After fitting, the
tree is flattened into arrays for vectorised prediction.
\emph{Ties:} we keep the first best split (lowest feature index, then
lowest threshold). scikit-learn scans features in a random order, which
explains most structural differences (\S\ref{sec:bench}).

\subsection{Support-vector machine (from scratch)}
\texttt{PegasosSVM} minimises (\ref{eq:svm}) with labels mapped to
$\pm1$. The bias is folded in as a constant feature, as the project
statement suggests. Each epoch visits a fresh random permutation in
mini-batches of $k=32$. At step $t$, with $A_t^+$ the batch rows whose
hinge is active ($y_i\mathbf w^\top\mathbf x_i<1$), the update is
\begin{equation}
  \mathbf w\leftarrow(1-\eta_t\lambda)\,\mathbf w+\tfrac{\eta_t}{k}\textstyle\sum_{i\in A_t^+}y_i\mathbf x_i,
\end{equation}
followed by the optional Pegasos projection onto
$\lVert\mathbf w\rVert\le1/\sqrt\lambda$, which tames the very large early
steps. The schedule is a hyper-parameter: $1/(\lambda t)$ (default),
$\eta_0/\sqrt t$, or constant. $J$ is evaluated on the full training set
after every epoch. Features are standardised with training means and
standard deviations, because the hinge loss is scale-sensitive.
\texttt{margin\_width}${}=2/\lVert\mathbf w\rVert$, and
\texttt{support\_vector\_mask} returns the training rows with
$y_i(\mathbf w^\top\mathbf x_i+b)\le1$.
\emph{Non-linear extension:} we chose \textbf{random Fourier features}
($D=1000$) followed by the same linear solver. Kernelised Pegasos needs
$O(n)$ kernel evaluations per step and stores up to $n$ coefficients. With
$n\approx40$k that is impractical, whereas RFF keeps each step $O(D)$. A
unit test confirms that $\mathbf z^\top\mathbf z$ approximates the RBF
kernel to within 0.05, and that RFF-Pegasos solves XOR, which the linear
model cannot.

\subsection{Bonus models}
\emph{Ridge} solves $(X^\top X+\lambda I)\mathbf w=X^\top\mathbf y$ on
centred data, leaving the intercept unpenalised. \emph{Random forest}:
30 of our trees, each on a bootstrap sample with a random 33\% of the
features considered at every split; the forest averages predictions or
probabilities.

\section{Experimental Setup}
\emph{Metrics:} for Task A, RMSE, MAE and $R^2$ on $\log(\text{price})$,
plus MAPE and $R^2$ in AZN. For Task B, precision, recall and F1 for the
premium class, ROC-AUC and PR-AUC (average precision), and the confusion
matrix. All metrics are our own NumPy code and are tested against
\texttt{sklearn.metrics}.
\emph{Protocol:} hyper-parameters are chosen on the validation split only.
Final models are fit on the training split and scored \emph{once} on the
test split, with a bootstrap 95\% CI (1000 resamples). For uncertainty we
also run 5-fold CV on train+val, stratified on price deciles, re-fitting
the featurizer, tier threshold and scaler inside every fold.
\emph{Search grids:} trees use \texttt{max\_depth} $\in\{2,\dots,25\}$
$\times$ \texttt{min\_samples\_leaf} $\in\{1,5,20,50\}$ (Task A) and
$\{1,20\}$ $\times$ \{Gini, entropy\} (Task B). The SVM uses
$\lambda\in[10^{-5},10^{-1}]$ with 30 epochs; when validation F1 is within
0.002, we take the larger $\lambda$, i.e.\ the wider margin. RFF uses
$\gamma\in[10^{-3},10^{-1}]\times\lambda\in\{10^{-3},10^{-4}\}$.
\emph{Baselines with matched settings:} \texttt{DecisionTreeRegressor} and
\texttt{DecisionTreeClassifier} with the same depth, leaf size and
criterion; \texttt{SGDClassifier}(hinge, $\alpha=\lambda$, same epochs);
\texttt{LinearSVC}($C=1/(\lambda n)$); \texttt{RBFSampler}+SGD with the
same $\gamma$, $D$ and $\lambda$; an RBF \texttt{SVC} on a 10k-row
subsample, since exact SVC is quadratic in $n$; and \texttt{Ridge} and
\texttt{RandomForestRegressor} with the same settings. Every scikit-learn
SVM is wrapped in a \texttt{Pipeline} with a \texttt{StandardScaler}.
Trivial mean and majority baselines anchor the scales. All randomness uses
seed 42. \texttt{python -m src.run\_all} regenerates every number, table
and figure in 11 minutes on a laptop CPU.

\section{Results}
\subsection{Tuning (validation only)}
\emph{Tree, Task A} (Fig.~\ref{fig:tree}a). With
\texttt{min\_samples\_leaf}${}=1$, training RMSE falls steadily with depth
while validation RMSE bottoms out at depth 10--12 (0.240) and then rises:
this is the variance side of the bias--variance trade-off. Raising the leaf
size to 5 or 20 removes the rise and lowers validation RMSE to 0.226. We
select depth 17 with leaf size 5. The learning curve (Fig.~\ref{fig:tree}c)
shows validation RMSE still falling from 0.308 with 2k rows to 0.226 with
40k rows, with a persistent train/validation gap, so more data would still
help. \emph{Tree, Task B} (Fig.~\ref{fig:tree}b). Gini beats entropy at
every depth up to 12. Both reach about 0.925 validation F1 around depth
14--17 and fit the training set perfectly by depth 25. Leaf size 20 costs
about 0.008 F1 here. We select Gini, depth 14, leaf size 1.

\begin{figure*}[t]
  \centering
  \includegraphics[width=0.49\linewidth]{tree_regression_tuning.pdf}\hfill
  \includegraphics[width=0.49\linewidth]{tree_classification_tuning.pdf}\\[2pt]
  \includegraphics[width=0.30\linewidth]{tree_learning_curve.pdf}\hfill
  \includegraphics[width=0.68\linewidth]{svm_lambda_sweep.pdf}
  \caption{Validation studies. (a)~Task A: depth sweep with train vs.\
  validation RMSE, and a second sweep over \texttt{min\_samples\_leaf}.
  (b)~Task B: Gini vs.\ entropy across depths and leaf sizes.
  (c)~Learning curve of the selected tree. (d)~SVM: as $\lambda$ grows,
  the margin $2/\lVert\mathbf w\rVert$ widens and more points become
  support vectors; validation F1 and AUC peak at
  $\lambda\approx10^{-3}$--$10^{-2}$ (dotted line: selected).}
  \label{fig:tree}
\end{figure*}

\emph{SVM} (Table~\ref{tab:lambda}, Fig.~\ref{fig:tree}d,
Fig.~\ref{fig:svm}). Increasing $\lambda$ from $10^{-5}$ to $10^{-1}$
shrinks $\lVert\mathbf w\rVert$ from 18.9 to 1.13. The margin therefore
widens 17-fold, and the share of support vectors rises from 16\% to 41\%.
Validation F1 is flat at 0.906--0.909 for
$\lambda\in[3\cdot10^{-4},3\cdot10^{-2}]$ and drops at both ends:
over-regularised at large $\lambda$, and not converged in 30 epochs at
$\lambda=10^{-5}$, where $J=0.80$. We select $\lambda=10^{-2}$.
Its objective falls from $J=1.00$ at $\mathbf w=\mathbf 0$ to 0.270 after
one epoch and 0.256 after 30. All three learning-rate schedules converge
to the same $J\approx0.256$, but the constant step is noisier. Running
more epochs widens the margin slightly, from 0.895 after one epoch to 0.960
after 40. The 2-D slice in Fig.~\ref{fig:svm} makes the effect of
$\lambda$ visible: $\lambda=1$ gives a margin band of width 3.8 that
contains most points. For RFF, the best setting is a wide kernel
($\gamma=0.003$, $\lambda=10^{-4}$); validation F1 falls steadily for
$\gamma\ge0.03$.

\begin{table}[t]
  \centering\small
  \caption{Pegasos $\lambda$ sweep (validation, 30 epochs).}
  \label{tab:lambda}
  \begin{tabular}{@{}lccccc@{}}
\toprule
$\lambda$ & val F1 & val AUC & $\|\mathbf{w}\|$ & $2/\|\mathbf{w}\|$ & SV \% \\
\midrule
0.1 & 0.897 & 0.965 & 1.13 & 1.771 & 41.4 \\
0.03 & 0.906 & 0.968 & 1.61 & 1.242 & 32.7 \\
0.01 & 0.908 & 0.969 & 2.09 & 0.959 & 27.7 \\
0.003 & 0.909 & 0.970 & 2.76 & 0.725 & 24.7 \\
0.001 & 0.909 & 0.970 & 3.71 & 0.539 & 23.3 \\
0.0003 & 0.907 & 0.968 & 4.55 & 0.440 & 22.2 \\
0.0001 & 0.894 & 0.961 & 5.68 & 0.352 & 20.4 \\
1e-05 & 0.868 & 0.938 & 18.85 & 0.106 & 16.0 \\
\bottomrule
\end{tabular}

\end{table}

\begin{figure*}[t]
  \centering
  \includegraphics[width=0.30\linewidth]{svm_objective.pdf}\hfill
  \includegraphics[width=0.68\linewidth]{svm_schedule_epochs.pdf}\\[2pt]
  \includegraphics[width=0.98\linewidth]{svm_boundary_2d.pdf}
  \caption{Pegasos diagnostics. Top: the training objective $J$ per epoch
  for several $\lambda$; the effect of the learning-rate schedule; and
  margin width vs.\ epochs. Bottom: a 2-D illustration on two standardised
  features (log area, distance to centre; 600 training rows). Solid line:
  decision boundary; dashed lines: $f(\mathbf x)=\pm1$; circled points:
  support vectors. Larger $\lambda$ gives a wider margin and more support
  vectors. The RFF panel shows a curved boundary.}
  \label{fig:svm}
\end{figure*}

\subsection{Test-set comparison}
\label{sec:bench}
Table~\ref{tab:results} summarises the headline results, and
Tables~\ref{tab:resA}--\ref{tab:resB} list every model. \emph{Task A:} our
tree reaches RMSE $0.232$ on $\log(\text{price})$ (95\% bootstrap CI
$[0.225, 0.240]$), $R^2=0.864$ and MAPE 15.4\%. In AZN this is an MAE of
44.6\,k and $R^2=0.763$, about 2.7 times better than predicting the mean.
\emph{Task B:} our tree reaches F1 $0.921$ (CI $[0.915, 0.928]$). The
linear SVM reaches F1 $0.900$ (CI $[0.893, 0.907]$) but a much higher
ROC-AUC (0.965 vs.\ 0.917), because the tree's scores are coarse leaf
frequencies while the SVM produces a continuous ranking. RFF matches but
does not beat the linear SVM. Once price, area and distance are on log or
standardised scales, the premium/standard boundary is close to linear, and
the trees' advantage comes from interactions that a wide RBF kernel
smooths away.

\begin{table}[t]
  \centering\small
  \caption{Headline results: test value (5-fold CV mean $\pm$ std).}
  \label{tab:results}
  \setlength{\tabcolsep}{3pt}
  \begin{tabular}{@{}lcc@{}}
\toprule
Model & Task A RMSE (log) & Task B F1 \\
\midrule
Ours: decision tree & 0.232 (0.232$\pm$0.004) & 0.921 (0.926$\pm$0.004) \\
Ours: SVM (Pegasos) & n/a & 0.900 (0.907$\pm$0.003) \\
Ours: SVM + RFF & n/a & 0.900 (0.906$\pm$0.003) \\
sk: decision tree & 0.232 (0.233$\pm$0.004) & 0.921 (0.926$\pm$0.004) \\
sk: SGD hinge & n/a & 0.900 (0.908$\pm$0.003) \\
\bottomrule
\end{tabular}

\end{table}

\begin{table*}[t]
  \centering\small
  \caption{Task A, regression on $\log(\text{price})$, test set. CV = 5-fold
  mean $\pm$ std on train+val (the forests were not cross-validated, to save
  time). $^b$bonus. Fit time on one CPU core.}
  \label{tab:resA}
  \begin{tabular}{@{}lcccccr@{}}
\toprule
Model & RMSE & MAE & $R^2$ & MAPE\% & CV RMSE & fit s \\
\midrule
Mean baseline & 0.629 & 0.487 & 0.000 & 55.2 & 0.630$\pm$0.003 & 0.00 \\
Ours: decision tree & 0.232 & 0.148 & 0.864 & 15.4 & 0.232$\pm$0.004 & 6.60 \\
sk: DecisionTreeRegressor & 0.232 & 0.148 & 0.864 & 15.4 & 0.233$\pm$0.004 & 0.53 \\
Ours: Ridge$^{b}$ & 0.242 & 0.168 & 0.852 & 17.1 & 0.240$\pm$0.004 & 0.03 \\
sk: Ridge & 0.242 & 0.168 & 0.852 & 17.1 & 0.240$\pm$0.004 & 0.08 \\
Ours: random forest$^{b}$ & 0.185 & 0.115 & 0.914 & 11.7 & --- & 131.76 \\
sk: RandomForestRegressor & 0.185 & 0.116 & 0.914 & 11.8 & --- & 4.36 \\
\bottomrule
\end{tabular}

  \vspace{4pt}
  \caption{Task B, premium vs.\ standard tier, test set (positive class =
  premium). The majority baseline predicts ``standard'' for every row.}
  \label{tab:resB}
  \begin{tabular}{@{}lcccccr@{}}
\toprule
Model & Prec. & Rec. & F1 & AUC & CV F1 & fit s \\
\midrule
Majority baseline & 0.000 & 0.000 & 0.000 & 0.500 & 0.000$\pm$0.000 & 0.00 \\
Ours: decision tree & 0.918 & 0.925 & 0.921 & 0.917 & 0.926$\pm$0.004 & 4.28 \\
sk: DecisionTreeClassifier & 0.919 & 0.923 & 0.921 & 0.916 & 0.926$\pm$0.004 & 0.47 \\
Ours: Pegasos SVM (linear) & 0.894 & 0.906 & 0.900 & 0.965 & 0.907$\pm$0.003 & 2.06 \\
sk: SGDClassifier hinge & 0.893 & 0.907 & 0.900 & 0.965 & 0.908$\pm$0.003 & 0.37 \\
sk: LinearSVC & 0.895 & 0.910 & 0.902 & 0.968 & 0.909$\pm$0.003 & 0.73 \\
Ours: RFF Pegasos SVM & 0.883 & 0.918 & 0.900 & 0.967 & 0.906$\pm$0.003 & 5.44 \\
sk: RBFSampler + SGD hinge & 0.886 & 0.915 & 0.901 & 0.967 & 0.907$\pm$0.003 & 2.23 \\
sk: SVC rbf (10k sub.) & 0.898 & 0.901 & 0.899 & 0.967 & --- & 1.93 \\
\bottomrule
\end{tabular}

\end{table*}

\emph{Ours vs.\ scikit-learn.} On a fully determined 8-row example, both
our classification trees (Gini and entropy) and our regression tree
reproduce scikit-learn's split features and thresholds exactly
(\texttt{tests/test\_decision\_tree.py}). On five random continuous
regression problems, the predictions agree to $10^{-10}$. On the real data
(Table~\ref{tab:tvs}) the trees have the same depth and leaf counts within
4, and agree on 94.5\% (A) and 98.9\% (B) of test predictions. Every metric
matches to within 0.002. The differences come from ties. In Task A the
root split is \texttt{area\_m2}${}\le100.1$ for us and
\texttt{log\_area}${}\le4.606$ for scikit-learn. Since $e^{4.606}=100.1$,
both produce exactly the same partition with the same gain: we keep the
lower feature index, while scikit-learn's random feature order picked the
other. Classification gains are ratios of small integers, so exact ties are
even more common deeper in the tree. Thresholds also differ in the last
float bits (90.15 vs.\ scikit-learn's float32 90.149998). Pegasos and
\texttt{SGDClassifier} agree to the third decimal on every metric, and
\texttt{LinearSVC}, which solves the dual exactly, is only 0.0025 F1 higher.
Our Ridge coefficients equal scikit-learn's to $8\cdot10^{-11}$.
\emph{Speed:} our tree fits in 6.6\,s (A) and 4.3\,s (B), against
scikit-learn's compiled Cython at 0.5\,s, i.e.\ 9--12$\times$ slower.
Pegasos takes 2.1\,s against 0.4\,s for SGD. Our 30-tree forest takes
132\,s against 4.4\,s, but is equally accurate (RMSE 0.185;
Fig.~\ref{fig:err}g). Prediction is under 0.01\,s for every model except
exact SVC, which needs 4.9\,s even when trained on only 10k rows; that cost
is why we did not use it as a main model.

\begin{table}[t]
  \centering\scriptsize
  \setlength{\tabcolsep}{2.5pt}
  \caption{Tree structure: ours vs.\ scikit-learn (same hyper-parameters).}
  \label{tab:tvs}
  \begin{tabular}{@{}lcccc@{}}
\toprule
Task & leaves & depth & root split (ours / sklearn) & same pred.\% \\
\midrule
A & 3945 / 3944 & 17 / 17 & area\_m2$\le$100.1 / log\_area$\le$4.606 & 94.5 \\
B & 1521 / 1525 & 14 / 14 & area\_m2$\le$90.15 / area\_m2$\le$90.15 & 98.9 \\
\bottomrule
\end{tabular}

\end{table}

\section{Error Analysis and Discussion}
\emph{What the models use.} Area dominates the tree importances (0.60 for
A, 0.54 for B), followed by distance to the centre (0.22, 0.15) and raw
lat/lng (Fig.~\ref{fig:imp}). The largest SVM weights tell a consistent
story: $+1.49$ for log area, $+0.55$ for rooms, and $+0.40$ for the
premium Ag \c{S}eher (``White City'') development; $-0.31$ for the Absheron district, $-0.27$
for houses/villas, and $-0.27$ for distance to the centre.

\emph{Luxury and bargain outliers.} The tree shrinks predictions toward
the middle (Fig.~\ref{fig:err}a). It over-predicts the cheapest decile by
$+0.11$ in log terms (12\%) and under-predicts the most expensive decile by
$-0.12$. MAPE is about 12\% in the middle deciles but 26\% and 22\% at the
two ends. Because leaves average, the most extreme prices cannot be
reached. The worst errors (\texttt{results/error\_worst.csv}) are
implausibly cheap listings, e.g.\ a 120\,m$^2$ new build near the
20~Yanvar metro listed at 47\,k~AZN. These are almost certainly down
payments or monthly instalments quoted as the price, which our price/m$^2$
floor of 300 does not catch.

\emph{Houses and sparse districts.} Houses/villas have 32\% MAPE against
12\% for apartments. Their value depends on land, build quality and plot
shape, and the scrape describes these poorly. Micro-locations with fewer
than 150 training listings have 29.9\% MAPE against 14.3\% for well-covered
ones, and test MAPE correlates with $\log(\text{train count})$ at $r=-0.53$
(Fig.~\ref{fig:err}c). The k-means clusters (bonus, Fig.~\ref{fig:clu})
show the same pattern geographically: RMSE is 0.19--0.21 in the dense
central and eastern Baku clusters (9--25k training rows) and 0.34--0.41 on
the northern Absheron coast. In clusters outside Absheron with only 2--26
training rows, RMSE reaches 0.85--0.89.

\emph{Tier errors sit at the boundary.} For listings within 5\% of the
threshold price, 38\% (tree) and 46\% (SVM) of predictions are wrong.
Beyond 50\% from the threshold, the error rates are 1.9\% and 1.4\%
(Fig.~\ref{fig:err}d). Most Task-B errors are therefore unavoidable
label noise: two near-identical flats a few thousand AZN apart fall on
opposite sides of an arbitrary median.

\begin{figure*}[t]
  \centering
  \includegraphics[width=0.49\linewidth]{error_by_price_decile.pdf}\hfill
  \includegraphics[width=0.245\linewidth]{error_by_location_frequency.pdf}\hfill
  \includegraphics[width=0.245\linewidth]{tier_error_vs_distance.pdf}\\[2pt]
  \includegraphics[width=0.30\linewidth]{roc_curves.pdf}\hfill
  \includegraphics[width=0.38\linewidth]{confusion_matrices.pdf}\hfill
  \includegraphics[width=0.29\linewidth]{bonus_random_forest.pdf}
  \caption{Error analysis on the test set. (a,b)~The tree shrinks extreme
  prices toward the middle; relative error is highest at both ends.
  (c)~Rarely listed micro-locations have higher error. (d)~Tier
  misclassifications concentrate near the threshold. (e)~ROC curves: the
  SVM ranks better (higher AUC) even though the tree has the higher F1.
  (f)~Confusion matrices. (g)~Bonus: the RMSE of our random forest falls
  from 0.264 with one tree to 0.185 with 30 trees.}
  \label{fig:err}
\end{figure*}

\begin{figure}[t]
  \centering
  \includegraphics[width=\linewidth]{feature_importance.pdf}\\[2pt]
  \includegraphics[width=0.75\linewidth]{svm_weights.pdf}
  \caption{Impurity-decrease importances of our trees (top) and the
  largest Pegasos weights on standardised features (bottom).}
  \label{fig:imp}
\end{figure}

\begin{figure}[t]
  \centering
  \includegraphics[width=0.95\linewidth]{bonus_clusters_pca.pdf}
  \caption{Bonus: k-means ($k=8$) on training coordinates, used to break
  down errors by area. A PCA projection of the test rows shows that the
  tiers overlap heavily in the first two components (7.7\% and 4.3\% of
  the variance), so no single direction separates them.}
  \label{fig:clu}
\end{figure}

\emph{Limitations.} (1)~These are \emph{asking} prices from an 11-day
snapshot; transaction prices and seasonal effects are not observed.
(2)~Stratifying the split on price deciles uses each row's own price to
assign it to a split. This never enters a feature, a threshold or a
model, but strict deployment would use a time-based split.
(3)~The fixed outlier rules remove some genuine luxury listings above
5\,M~AZN, and keep some mis-entered instalment prices.
(4)~The free-text description, which mentions renovations, documents and
views, is unused.
(5)~The forests and the exact SVC were evaluated only on the test split,
not with CV.
(6)~Our NumPy tree is about 10$\times$ slower than compiled code, which
limited how many forest configurations we could tune.

\section{Conclusion}
A carefully cleaned and leakage-free version of the bina.az scrape lets
simple from-scratch models predict Baku sale prices to within about 15\%
MAPE and classify the price tier with F1 0.92. Our NumPy tree and Pegasos
SVM match scikit-learn to within 0.002 on every metric; the remaining
differences come from tie-breaking and float precision, not from
modelling. The decision tree is the stronger classifier at a fixed
threshold, while the SVM is the better ranker. Averaging our own trees in a
forest gives the largest single gain (RMSE $0.232\rightarrow0.185$). Next
we would add text features from the descriptions, use a time-based split,
target-encode the sparse micro-locations, and try gradient boosting with
our trees.

\section*{Tools and Acknowledgements}
\textbf{AI-assistant disclosure.} AI coding assistants (Claude Code, Anthropic, ChatGPT) were used to draft the structure of the code(its pipeline). The idea of tests, figures and the first version of the report text were proposed by assistants. \textbf{Data:} bina.az sale
listings, Kaggle dataset \url{sehriyarmemmedli/binaaz-sale-project}; we
do not redistribute the raw data or personal fields. \textbf{Software:}
NumPy, pandas and Matplotlib; scikit-learn was used only for baselines,
bonus clustering/PCA, and sanity checks in tests.

\begin{thebibliography}{9}
\bibitem{pegasos}
S.~Shalev-Shwartz, Y.~Singer, N.~Srebro, and A.~Cotter,
``Pegasos: primal estimated sub-gradient solver for SVM,''
\emph{Mathematical Programming}, vol.~127, no.~1, pp.~3--30, 2011.
\bibitem{cart}
L.~Breiman, J.~Friedman, R.~Olshen, and C.~Stone,
\emph{Classification and Regression Trees}. Wadsworth, 1984.
\bibitem{rff}
A.~Rahimi and B.~Recht, ``Random features for large-scale kernel
machines,'' in \emph{Advances in Neural Information Processing Systems
(NIPS)}, 2007, pp.~1177--1184.
\bibitem{sklearn}
F.~Pedregosa \emph{et al.}, ``Scikit-learn: Machine learning in Python,''
\emph{J. Mach. Learn. Res.}, vol.~12, pp.~2825--2830, 2011.
\end{thebibliography}

\end{document}
